\section{本章小结}
\label{sec:graph-walk-summary}

本章以图上的随机游走为主线，将前几章建立的矩阵幂、归一化拉普拉斯谱、Cheeger导通率以及电网络有效电阻融会贯通。

在位置分布与混合速度方面，随机游走的一步转移规则由行随机矩阵$\bm P=\bm D^{-1}\bm W$描述，多步分布由$\bm\nu^{(t)}=\bm\nu^{(0)}\bm P^t$给出。无向连通图的详细平衡与拉普拉斯零空间保证了按度数分配的平稳分布$\pi_i=d_i/\operatorname{vol}(V)$唯一存在。

通过度平方根相似变换$\bm S=\bm D^{1/2}\bm P\bm D^{-1/2}=\bm I-\mathcal L$，从任意初始分布出发均收敛于$\bm\pi$当且仅当图不是二部图（$\theta_n>-1$）；引入惰性游走$\bm Q=(\bm I+\bm P)/2$可消除负特征值奇偶振荡，使非平稳谱模长由$\gamma=\mu_2/2$决定，具体误差界还含初态因子，并经由Cheeger不等式$\phi_*^2/4\leq\gamma\leq\phi_*$与图的瓶颈导通率$\phi_*$紧密相连。

在首达时间方面，不同起终点间的最大期望在完全图（$\Theta(n)$）、路径与环（$\Theta(n^2)$）、超立方体（$\Theta(n)$）以及棒棒糖图（$\Theta(n^3)$）中展示了不同的等待时间量级。一般图上首达期望满足接地拉普拉斯方程$\bm L_{-j,-j}\bm h_{-j,j}=\bm d_{-j}$，而两点间的往返时间则满足电学恒等式$h_{i,j}+h_{j,i}=\operatorname{vol}(V)R_{i,j}$，对连通简单单位权图，由此得到首达时间上界$2m\,\operatorname{dist}(i,j)\leq 2m(n-1)<n^3$。

进一步，Matthews定理通过调和数$H_{n-1}$将遍历时间$C$夹在$h_{\min}H_{k-1}$与$h_{\max}H_{n-1}$之间，而生成树深度优先巡游则利用有效电阻给出了消除多余对数因子的遍历上界$C\leq\operatorname{vol}(V)\sum_{e\in T}R_e$，在单位权图上进一步有$C\leq2m(n-1)$。

在有向图PageRank中，死胡同修补与阻尼跳转给出唯一平稳分布；只有跳转分布处处为正时，才直接保证所有顶点得分严格为正。非平稳特征值模长不超过阻尼系数$\alpha$，稀疏幂法具有几何收缩界与残差误差证书。

个性化PageRank一节改用$\alpha$表示重启概率：局部推送在不变量$\mathcal R(\bm s)=\bm p+\mathcal R(\bm r)$下具有$O(1/(\alpha\varepsilon))$的推送工作量和单位度数误差界，另需种子初始化成本。它是否为亚线性算法取决于参数；局部切分还需排序、按原图计算割权，以及额外的种子与质量保证条件。
